% Generated from locale/tr/content/first-order-logic/syntax-and-semantics/main-operator.tex; do not hand-edit.


\olfileid[tr]{fol}{syn}{mai}

\olsection{\printtoken{S}{main operator} Kavramı}

\begin{explain}
Bir !!a{formula}~$!A$ kurulurken en son kullanılan işleçten söz etmek
çoğu kez yararlıdır. Bu işlece $!A$'nın \emph{ana işleci} denir.
Sezgisel olarak bu, $!A$'nın ``en dıştaki'' işlecidir. Örneğin
$\lnot !A$'nın ana işleci $\lnot$, $(!A \lor !B)$'nin ana işleci
$\lor$'dur vb.
\end{explain}


\begin{defn}[!!^{main operator}]
\ollabel{def:main-op}
\emph{!!{main operator}}, bir !!a{formula}~$!A$ için aşağıdaki gibi
tanımlanır:
\begin{enumerate}
\item \indcase*{!A}{!A}{$\indfrm$ için !!{main operator} yoktur.}

\tagitem{prvNot}{\indcase{!A}{\lnot !B}{$\indfrm$ için !!{main operator},
  $\lnot$'tur.}}{}

\tagitem{prvAnd}{\indcase{!A}{(!B \land !C)}{$\indfrm$ için
  !!{main operator}, $\land$'dır.}}{}

\tagitem{prvOr}{\indcase{!A}{(!B \lor !C)}{$\indfrm$ için
  !!{main operator}, $\lor$'dur.}}{}

\tagitem{prvIf}{\indcase{!A}{(!B \lif !C)}{$\indfrm$ için
  !!{main operator}, $\lif$'tir.}}{}

\tagitem{prvIff}{\indcase{!A}{(!B \liff !C)}{$\indfrm$ için
  !!{main operator}, $\liff$'tir.}}{}

\tagitem{prvAll}{\indcase{!A}{\lforall[x][!B]}{$\indfrm$ için
  !!{main operator}, $\lforall$'dır.}}{}

\tagitem{prvEx}{\indcase{!A}{\lexists[x][!B]}{$\indfrm$ için
  !!{main operator}, $\lexists$'tir.}}{}
\end{enumerate}
\end{defn}

Her durumda formülde özellikle gösterilen !!{main operator}
\emph{oluşumunu} kastediyoruz. Örneğin $((!D \lif !E) \lif (!E \lif
!D))$ formülü, $!B$'nin $(!D \lif !E)$ ve $!C$'nin $(!E \lif !D)$
olduğu $(!B \lif !C)$ biçiminde olduğundan, $\lif$'in ikinci oluşumu
!!{main operator}{}tir.

\begin{explain}
Bu, atomik olmayan bütün !!{formula}s{}in her biri için !!{main operator}
oluşumunu veren bir fonksiyonun \emph{özyinelemeli} tanımıdır.
!!{formula}s tümevarımsal olarak tanımlandığı için her
!!{formula}~$!A$, \olref{def:main-op}'taki durumlardan birini sağlar.
Bu, atomik olmayan her !!{formula}~$!A$ için bir !!a{main
  operator} bulunduğunu güvence altına alır. Her !!{formula} bu
koşullardan yalnızca birini sağladığı ve $!A$'nın her durumda kendilerinden
kurulduğu daha küçük !!{formula}s tek türlü belirlendiği için $!A$'nın
!!{main
  operator} oluşumu tektir; dolayısıyla bir fonksiyon tanımlamış oluruz.
\end{explain}

!!{formula}s, sahip oldukları !!{main operator} simgesine göre
\olref{tab:main-op}'taki adlarla anılır.
\iftag{defNot,defOr,defAnd,defIf,defIff,defTrue,defFalse,defEx,defAll}
{ Ancak tanımlı işleçlerin resmen !!{formula}s içinde yer almadığını
anımsayın. Bunlar yalnızca kısaltmadır; dolayısıyla resmen bir formülün
ana işleci olamazlar. Bu yüzden bütün !!{formula}s hakkındaki ispatlarda
ayrı olarak ele alınmaları gerekmez.}

\begin{table}[!h]
\centering
\begin{tabular}{c | c | c}
!!^{main operator} & !!{formula} türü & Örnek\\
\hline
yok & atomik (!!{formula}) &
\iftag{prvFalse}{$\lfalse$,}{}
\iftag{prvTrue}{$\ltrue$,}{}
$\Atom{R}{t_1, \dots, t_n}$,
$\eq[t_1][t_2]$\\
$\lnot$ & değilleme & $\lnot !A$ \\
$\land$ & ve bağlacı & $(!A \land !B)$ \\
$\lor$ & veya bağlacı & $(!A \lor !B)$ \\
$\lif$ & !!{conditional} & $(!A \lif !B)$ \\
$\liff$ & !!{biconditional} & $(!A \liff !B)$ \\
$\lforall[][]$ & evrensel (!!{formula})& $\lforall[x][!A]$ \\
$\lexists[][]$ & varlıksal (!!{formula})& $\lexists[x][!A]$
\end{tabular}
\caption{Ana işleçlere göre !!{formula}s ve adları}
\ollabel{tab:main-op}
\end{table}

